\documentclass[a4paper,12pt]{article}
\usepackage[T2A]{fontenc}
\usepackage[utf8]{inputenc}
\usepackage[russian]{babel}
\usepackage{amsmath, amssymb, amsthm}
\usepackage{amsfonts}
\usepackage{geometry}
\geometry{left=2cm,right=2cm,top=2cm,bottom=2cm}

\title{Текст речи к защите магистерской диссертации}
\author{Зайков Иван Викторович}
\date{2026}

\begin{document}

\maketitle

\section*{Слайд 1. Титульный}

Уважаемый председатель, уважаемые члены Государственной экзаменационной комиссии!

Вашему вниманию представляется магистерская диссертация на тему: \textbf{«Задача о минимальной норме интерполяционного проектора»}.

Работа выполнена на кафедре математического анализа Ярославского государственного университета имени П. Г. Демидова под руководством кандидата физико-математических наук, доцента Ухалова Алексея Юрьевича.

\section*{Слайд 2. Постановка задачи}

В работе рассматривается задача линейной интерполяции непрерывной функции на единичном кубе. Для функции $f(x)$, заданной на кубе $Q_n$, строится линейный полином
\[
p(x) = a_1 x_1 + \dots + a_n x_n + b,
\]
который совпадает с функцией в $n+1$ узле — вершинах невырожденного симплекса $S$. Это отображение называется интерполяционным проектором.

Норма проектора, или константа Лебега, вычисляется как максимум суммы модулей базисных многочленов Лагранжа по всем вершинам куба:
\[
\|P\| = \max_{x \in \operatorname{ver}(Q_n)} \sum_{j=1}^{n+1} |\lambda_j(x)|.
\]
Это связано с тем, что функция $L(x) = \sum |\lambda_j(x)|$ является выпуклой и кусочно-линейной на $Q_n$, поэтому её максимум достигается в одной из вершин куба.

Особое значение имеет минимальная норма интерполяционного проектора
\[
\theta_n = \min_{x^{(j)} \in Q_n} \|P\|,
\]
которая достигается при оптимальном выборе узлов интерполяции (вершин симплекса). Задача нахождения $\theta_n$ является классической в теории приближений.

Цель нашей работы — получить верхнюю оценку $\theta_{31}$, используя симплекс максимального объёма, построенный на основе матрицы Адамара.

\section*{Слайд 3. Обозначения}

Введём основные обозначения. $Q_n = [0,1]^n$ — единичный куб в $n$-мерном пространстве. $C(Q_n)$ — пространство непрерывных функций, определённых на $Q_n$, с равномерной нормой
\[
\|f\|_{C(Q_n)} = \max_{x \in Q_n} |f(x)|.
\]

$\Pi_1(\mathbb{R}^n) = \operatorname{span}\{1, x_1, \dots, x_n\}$ — линейная оболочка функций $1, x_1, \dots, x_n$, то есть пространство всех линейных функций от $n$ переменных.

Симплекс $S \subset Q_n$ — это невырожденный многогранник, вершины которого служат узлами интерполяции. Невырожденность означает, что его вершины не лежат в одной гиперплоскости (эквивалентно, $\det A \neq 0$), что гарантирует единственность интерполяционного полинома.

\section*{Слайд 4. Базисные многочлены Лагранжа}

Для интерполяции функции в вершинах симплекса удобно использовать базисные многочлены Лагранжа $\lambda_j(x)$. Они строятся через барицентрические координаты относительно вершин симплекса.

Матрица $A$ составлена из координат вершин симплекса, расположенных в строках, и дополненных единицей в последнем столбце:
\[
A = \begin{pmatrix}
x_1^{(1)} & \dots & x_n^{(1)} & 1 \\
\vdots & & \vdots & \vdots \\
x_1^{(n+1)} & \dots & x_n^{(n+1)} & 1
\end{pmatrix}, \qquad \det A \neq 0.
\]

Базисные многочлены Лагранжа определяются как
\[
\lambda_j(x) = \frac{\Delta_j(x)}{\Delta},
\]
где $\Delta = \det(A)$, а $\Delta_j(x)$ — определитель, полученный заменой $j$-й строки на $(x_1,\dots,x_n,1)$. Они удовлетворяют главному свойству:
\[
\lambda_j(x^{(k)}) = \delta_{jk}, \qquad \sum_{j=1}^{n+1} \lambda_j(x) = 1.
\]

На практике удобнее вычислять их через обратную матрицу:
\[
\lambda(x) = A^{-1} b(x),
\]
где $b(x) = (x_1, \dots, x_n, 1)^T$ — расширенный вектор точки.

\section*{Слайд 5. Интерполяционный проектор}

Дадим строгое определение проектора. Это отображение $P: C(Q_n) \to \Pi_1(\mathbb{R}^n)$, которое каждой функции $f$ ставит в соответствие интерполяционный полином, удовлетворяющий условиям
\[
P f(x^{(j)}) = f(x^{(j)}), \quad j = 1,\dots,n+1.
\]

Интерполяционная формула Лагранжа имеет вид:
\[
P f(x) = \sum_{j=1}^{n+1} f(x^{(j)}) \lambda_j(x).
\]

Норма проектора вычисляется по формуле:
\[
\|P\| = \max_{x \in \operatorname{ver}(Q_n)} \sum_{j=1}^{n+1} |\lambda_j(x)|.
\]

Минимальная норма обозначается $\theta_n = \min \|P\|$, и наша задача — найти её верхнюю оценку.

\section*{Слайд 6. Неравенство Лебега}

Важную роль в теории приближений играет неравенство Лебега. Пусть $E_n(f)$ — наилучшее приближение функции $f$ полиномами из $\Pi_1(\mathbb{R}^n)$:
\[
E_n(f) = \min_{p \in \Pi_1(\mathbb{R}^n)} \|f - p\|_{C(Q_n)}.
\]

Тогда справедливо неравенство:
\[
\|f - P f\|_{C(Q_n)} \le (1 + \|P\|) \cdot E_n(f).
\]

Множитель $(1 + \|P\|)$ показывает, во сколько раз интерполяция может быть хуже наилучшего приближения. Чем меньше норма проектора, тем ближе интерполяционный полином к оптимальному.

\section*{Слайд 7. Общие оценки $\theta_n$}

Для величины $\theta_n$ известны общие оценки:
\[
\frac{1}{4}\sqrt{n} < \theta_n < 3\sqrt{n}, \qquad 3 - \frac{4}{n+1} \le \theta_n.
\]

Кроме того, существует асимптотическая эквивалентность:
\[
\|P\| \asymp \theta_n,
\]
что означает, что существуют константы $c_1, c_2 > 0$ такие, что
\[
c_1\sqrt{n} \le \theta_n \le c_2\sqrt{n}.
\]

Это означает, что для получения верхней оценки $\theta_n$ достаточно рассмотреть симплексы максимального объёма.

\section*{Слайд 8. Матрица Адамара}

Переходим к ключевому инструменту нашей работы — матрицам Адамара. Это квадратные матрицы из элементов $\pm 1$, строки которых попарно ортогональны:
\[
H H^T = n I_n.
\]

Их определитель максимален и равен $|\det H| = n^{n/2}$.

Матрицы Адамара существуют не для всех порядков: необходимо, чтобы порядок был равен $1$, $2$ или кратен $4$. Гипотеза Адамара утверждает, что они существуют для любого порядка, кратного $4$, но это пока не доказано. Для порядка $32$ известно $13\,710\,027$ неэквивалентных матриц.

Матрицы Адамара можно строить рекурсивно методом Сильвестра:
\[
H_{1} = (1), \qquad H_{2n} = \begin{pmatrix} H_n & H_n \\ H_n & -H_n \end{pmatrix}.
\]

Связь с нашей задачей: из матрицы Адамара порядка $n+1$ можно получить симплекс максимального объёма в $n$-мерном кубе.

\section*{Слайд 9. Типы матриц Адамара порядка 32}

Для исследования мы взяли шесть матриц порядка 32 из библиотеки Нила Слоуна:
\begin{itemize}
    \item \texttt{syl} — метод Сильвестра;
    \item \texttt{pal} — метод Пэли, построенный на основе конечных полей;
    \item \texttt{t1}, \texttt{t2}, \texttt{t3}, \texttt{t4} — метод тензорного произведения.
\end{itemize}

\section*{Слайд 10. Количество матриц Адамара}

Таблица показывает, как быстро растёт число неэквивалентных матриц Адамара с ростом порядка. Для порядка $32$ их уже более $13$ миллионов. Именно поэтому полный перебор всех симплексов невозможен, и мы ограничились шестью представителями различных типов.

\section*{Слайд 11. Построение симплекса}

Построение симплекса из матрицы Адамара выполняется в несколько шагов.

Исходные данные: матрица Адамара $H$ порядка $32$ с элементами $\pm 1$.

\textbf{Шаги:}
\begin{enumerate}
    \item Приводим $H$ к виду $B$, где первый столбец и первая строка состоят из единиц (умножением строк и столбцов на $-1$).
    \item Удаляем первую строку и первый столбец. Получаем матрицу $C$ порядка $31$.
    \item Заменяем элементы: $-1 \to 1$, $1 \to 0$. Получаем $(0/1)$-матрицу $D$ порядка $31$.
    \item Добавляем нулевую строку (вершину). Получаем $32$ вершины в $\mathbb{R}^{31}$.
\end{enumerate}

Результат:
\[
V = \{0, \operatorname{row}_1(D), \dots, \operatorname{row}_{31}(D)\} \subset [0,1]^{31}
\]
— симплекс максимального объёма.

\section*{Слайд 12. Программная реализация}

Для вычислений была разработана программа на C++ с использованием библиотеки Eigen. Ключевые классы:
\begin{itemize}
    \item \texttt{ISimplexProvider} — интерфейс для получения вершин симплекса;
    \item \texttt{SilvesterMethodProvider}, \texttt{HadamardFileProvider} — реализация;
    \item \texttt{FullEnumUpperBoundCalculator} — многопоточный перебор вершин;
    \item \texttt{LebesgueFunction} — вычисление суммы модулей;
    \item \texttt{HadamardMatrixIterator} — чтение матриц из файла.
\end{itemize}

Технологии: C++17, Eigen, CMake, многопоточность с атомарными операциями.

\section*{Слайд 13. Ввод и вывод данных}

Программа поддерживает два режима работы:
\begin{enumerate}
    \item Построение симплекса по Сильвестру;
    \item Загрузка матриц Адамара из файла.
\end{enumerate}

Управление осуществляется через параметры командной строки:
\begin{itemize}
    \item \texttt{--dimension N} — размерность пространства $\mathbb{R}^N$;
    \item \texttt{--load-hadamard FILE} — загрузить матрицу Адамара из файла;
    \item \texttt{--threads N} — количество потоков;
    \item \texttt{--range START END} — частичный перебор вершин в диапазоне.
\end{itemize}

\section*{Слайд 14. Полный перебор}

Гиперкуб в размерности $31$ содержит
\[
N = 2^{31} = 2\,147\,483\,648
\]
вершин.

Для их перебора мы используем:
\begin{itemize}
    \item кодирование вершины 31-битовым целым числом;
    \item многопоточность (16 потоков);
    \item атомарные операции для синхронизации;
    \item логирование каждые $10^7$ вершин.
\end{itemize}

Это позволило завершить перебор за $290$ секунд (около $5$ минут).

\section*{Слайд 15. Параметры вычислительной системы}

Вычисления проводились на следующем оборудовании:
\begin{itemize}
    \item Процессор: Intel Core i5-1240P (16 потоков);
    \item Оперативная память: 15 ГБ;
    \item Операционная система: Ubuntu 24.04.3 LTS;
    \item Компилятор: GCC 13.3.0;
    \item Библиотека: Eigen 3.
\end{itemize}

\section*{Слайд 16. Результаты}

Главный результат работы — полученная верхняя оценка:
\[
\boxed{\theta_{31} \le 5.000}.
\]

Теоретическая оценка составляет:
\[
\sqrt{32} \approx 5.656854.
\]

Таким образом, нам удалось уточнить существующую границу. Перебор всех вершин занял $290$ секунд при скорости около $7.4 \times 10^6$ вершин в секунду.

\section*{Слайд 17. Заключение}

В ходе работы выполнено:
\begin{enumerate}
    \item Построен симплекс максимального объёма в $\mathbb{R}^{31}$ на основе матрицы Адамара;
    \item Реализован многопоточный полный перебор $2^{31}$ вершин гиперкуба;
    \item Получена верхняя оценка $\theta_{31} \le 5.000$.
\end{enumerate}

\section*{Слайд 18. Спасибо за внимание}

Доклад окончен. Спасибо за внимание! Я готов ответить на ваши вопросы.

\end{document}